\section{The full generator cokernel map and the original zeta
divisor}\label{the-full-generator-cokernel-map-and-the-original-zeta-divisor}

24 September 2026. Proof locators GCF0--GCF8. This derives an actual map
in the existing coefficient sequence and its global comparison with the
original meromorphic zeta function. It does not establish a Deligne
weight bound.

\subsection{GCF0. Prerequisites and the precise receiving
object}\label{gcf0.-prerequisites-and-the-precise-receiving-object}

The supporting datum remains the user's \(\tau\), with the
information-layer notation \(Z_0,Z_1,Z_2,\ldots\). All operations in
this note occur after the complete arithmetic reconstruction. No
addition, arithmetic value, distance, or parity is assigned to \(\tau\).
Distinct reconstructed branches are not pooled.

The actual source comparison is SSI0--SSI8 in
\nolinkurl{EXACT_SCHWARTZ_SUMMATION_IMAGE.md}, together with S1--S7 in
\texttt{GLOBAL\_MELLIN\_SYNTHESIS.md}. It gives the topological
isomorphism \[
\Theta:A\longrightarrow\mathcal B,\qquad
\Theta a(s)=\frac12\int_0^\infty a(u)u^s\frac{du}{u},
\quad J=\Sigma\mathcal S_0^{\mathrm{even}},\quad
\Theta J=\mathcal I,\quad Q=A/J\simeq\mathcal Q=\mathcal B/\mathcal I.
\tag{GCF0.1}
\] Here \(A\) has the seminorms
\(\sup_{u>0}(u^N+u^{-N})|(u\partial_u)^ja(u)|\), and
\(\Sigma f(u)=2\sum_{k\ge1}f(ku)\). The coefficient space and its closed
subspace are \[
\mathcal B=\left\{F\text{ entire}: b_{R,M}(F):=
\sup_{|\Re s|\le R}(1+|\Im s|)^M|F(s)|<\infty
\text{ for every }R,M\in\mathbb N\right\},
\] \[
\mathcal I=\{F\in\mathcal B:F^{(j)}(\rho)=0
\text{ for every actual nontrivial zero }\rho\text{ of }\zeta,
\ 0\le j<m_\rho\}.
\tag{GCF0.2}
\] All multiplicities \(m_\rho\) are retained. The entire source
element, obtained from the original Gaussian-polynomial source, is \[
F_0(s)=c(s)\zeta(s),\qquad
c(s)=\frac{s(s-1)}8\pi^{-s/2}\Gamma(s/2).
\tag{GCF0.3}
\] It belongs to \(\mathcal I\), has exactly its prescribed zero
divisor, and satisfies \(F_0(0)=F_0(1)=1/8\). These are proved source
properties, not a definition of the original \(\zeta\). In particular
the recovered original series still has unit term \(1\) and every
integer term, and its Euler expansion still contains every prime-power
repetition. The multiplier in (GCF0.3) is never suppressed below.

The operator is the continuous map \(LF(s)=sF(s)\) on \(\mathcal B\),
restricted to \(\mathcal I\) and induced on \(\mathcal Q\). Under
(GCF0.1) it corresponds to \(-u\partial_u\). Thus neither a numerical
coordinate for \(\tau\) nor an unconstructed arithmetic action is being
introduced. CW6--CW8 already calculate its quotient characters and
resolvent; VSD13 proves its closed codimension-one images on
\(\mathcal B\) and \(\mathcal I\). RZ1--RZ6 retain the full quotient
resolvent and multiplicity projectors. The present new receiver is the
map between those two cokernels as the parameter varies, including its
original meromorphic divisor.

\subsection{GCF1. Both cokernel lines, with their actual quotient
maps}\label{gcf1.-both-cokernel-lines-with-their-actual-quotient-maps}

For every \(\lambda\in\mathbb C\), define \[
\epsilon_\lambda(F)=F(\lambda)\quad(F\in\mathcal B),\qquad
\ell_\lambda(F)=(F/F_0)(\lambda)\quad(F\in\mathcal I).
\tag{GCF1.1}
\] The second quotient is entire because the complete zero divisor of
\(F_0\) is imposed on every \(F\in\mathcal I\). It need not belong to
\(\mathcal B\). If \(F_0\) has order \(m\) at \(\lambda\), its removable
value is \(F^{(m)}(\lambda)/F_0^{(m)}(\lambda)\); the nonzero
denominator is retained.

Both maps are continuous and onto. For \(\epsilon_\lambda\), a section
is \(z\mapsto z e^{(s-\lambda)^2}\). For \(\ell_\lambda\), a section is
\(z\mapsto zF_0\). Continuity of the latter follows either by the
displayed derivative formula or by a contour estimate: on a circle
avoiding the zeros of \(F_0\), division by \(F_0\) is bounded, and the
maximum principle bounds the entire quotient on the enclosed compact set
by this boundary bound. A larger-strip seminorm controls the numerator
on that circle.

Division by \(s-\lambda\) is continuous on the closed subspace of
\(\mathcal B\) vanishing at \(\lambda\). Outside a fixed disk the
denominator is bounded away from zero. Inside the disk, the maximum
principle for the removable quotient bounds it on a larger circle. This
proves a finite larger-strip seminorm bound for every target seminorm.
On \(\mathcal I\), division preserves its prescribed orders precisely
when \(\ell_\lambda(F)=0\). Consequently the two exact rows are \[
0\to\mathcal B\xrightarrow{L-\lambda}\mathcal B
\xrightarrow{\epsilon_\lambda}\mathbb C\to0,
\qquad
0\to\mathcal I\xrightarrow{L-\lambda}\mathcal I
\xrightarrow{\ell_\lambda}\mathbb C\to0.
\tag{GCF1.2}
\] Every image is closed, and the inverse on the image is continuous.
The inclusion \(\mathcal I\hookrightarrow\mathcal B\) induces the exact
scalar map \[
\begin{array}{ccc}
\mathcal I/(L-\lambda)\mathcal I&\longrightarrow&
\mathcal B/(L-\lambda)\mathcal B\\
\ell_\lambda\downarrow\ \cong&&\cong\ \downarrow\epsilon_\lambda\\
\mathbb C&\xrightarrow{\ \cdot F_0(\lambda)\ }&\mathbb C.
\end{array}
\tag{GCF1.3}
\] Indeed \(F(\lambda)=F_0(\lambda)(F/F_0)(\lambda)\), including
removable values at every zero. Thus the full source expression in
(GCF0.3) is an actual cokernel comparison, not a chosen positivity
factor.

\subsection{GCF2. Holomorphic parameter
dependence}\label{gcf2.-holomorphic-parameter-dependence}

The quotient-line comparison is a holomorphic family over the entire
parameter plane, including the zero set. Here is an explicit
verification avoiding any assumption about holomorphic lifts through
\(\mathcal B\to\mathcal Q\).

Put \[
v_{\mathcal B,\lambda}(s)=e^{(s-\lambda)^2},\qquad
v_{\mathcal I,\lambda}(s)=F_0(s),\qquad
\eta_{\mathcal B,\lambda}=\epsilon_\lambda,\quad
\eta_{\mathcal I,\lambda}=\ell_\lambda.
\tag{GCF2.1}
\] For \(E=\mathcal B,\mathcal I\), the maps \[
S_{E,\lambda}F(s)=
\frac{F(s)-v_{E,\lambda}(s)\eta_{E,\lambda}(F)}{s-\lambda}
\tag{GCF2.2}
\] take values in \(E\), and direct multiplication gives \[
(L-\lambda)S_{E,\lambda}F
=F-v_{E,\lambda}\eta_{E,\lambda}(F),
\qquad S_{E,\lambda}(L-\lambda)F=F.
\tag{GCF2.3}
\] For \(E=\mathcal I\), the numerator has one additional order at
\(s=\lambda\) relative to the required divisor, since its quotient by
\(F_0\) vanishes there. This proves the asserted target of (GCF2.2),
also at a multiple zero.

On any compact set of parameters \(\lambda\), the division bounds in
GCF1 can use one larger disk and strip. The Gaussian and all its
parameter derivatives have uniform bounds in every \(\mathcal B\)
seminorm on that parameter set. The map \(F\mapsto F/F_0\) is continuous
from \(\mathcal I\) to entire functions with compact convergence, by the
same boundary-circle argument. Evaluations and all their parameter
derivatives are therefore bounded by finitely many source seminorms
uniformly on a compact parameter set. Formula (GCF2.2) has a removable
singularity along \(s=\lambda\); Cauchy's integral formula and these
bounds prove holomorphic dependence with values in every coefficient
seminorm. They also prove that (GCF2.2) carries a holomorphic
\(E\)-valued family to a holomorphic \(E\)-valued family.

Let \(\mathcal H_E(U)\) denote holomorphic \(E\)-valued maps on an open
parameter set \(U\subset\mathbb C\). Equations (GCF2.1)--(GCF2.3) prove,
on every such set, \[
\mathcal H_E(U)/(L-\lambda)\mathcal H_E(U)
\xrightarrow{\ \eta_{E,\lambda}\ }\mathcal O(U).
\tag{GCF2.4}
\] This is an isomorphism: the stated section proves surjectivity and
\(S_{E,\lambda}\) proves its exact kernel. Inclusion of the coefficient
spaces is consequently represented by the global map of these two
quotient line sheaves \[
\mathcal O_{\mathbb C}\xrightarrow{\ \cdot F_0\ }\mathcal O_{\mathbb C}.
\tag{GCF2.5}
\] Its determinant is the displayed one-by-one entry \(F_0\), with its
exact source factor \(1/8\). No determinant or trace of an
infinite-dimensional operator has been assumed. The frame of the first
line uses the actual source element \(F_0\); changing that frame would
change the written entry and its inverse frame map together.

\subsection{GCF3. Exact relation with the actual quotient
obstruction}\label{gcf3.-exact-relation-with-the-actual-quotient-obstruction}

Apply the kernel/cokernel sequence to multiplication by \(L-\lambda\) in
the actual row \(0\to\mathcal I\to\mathcal B\to\mathcal Q\to0\). Since
the first two multiplications are injective, it gives \[
0\to\ker(L-\lambda\mid\mathcal Q)
\xrightarrow{\partial_\lambda}\mathbb C
\xrightarrow{\ \cdot F_0(\lambda)\ }\mathbb C
\to\operatorname{coker}(L-\lambda\mid\mathcal Q)\to0.
\tag{GCF3.1}
\] The boundary has a literal formula. Represent a kernel class by
\(G\in\mathcal B\), so \((s-\lambda)G\in\mathcal I\). Then \[
\partial_\lambda[G]
=\left(\frac{(s-\lambda)G(s)}{F_0(s)}\right)_{s=\lambda}.
\tag{GCF3.2}
\] Changing \(G\) by an element of \(\mathcal I\) changes the numerator
quotient by a multiple of \(s-\lambda\), so this is well defined. These
computations prove (GCF3.1), including its sign, directly.

When \(F_0(\lambda)\ne0\), both outside terms vanish. When
\(\lambda=\rho\) is an actual nontrivial zero, both outside terms are
one-dimensional and \[
G_\rho(s)=\frac{F_0(s)}{s-\rho}\in\mathcal B,
\qquad \partial_\rho[G_\rho]=1.
\tag{GCF3.3}
\] This remains true for a multiple zero: the class has its last
retained nonzero jet of order \(m_\rho-1\). All statements quantify over
the full actual zero set.

This boundary is the obstruction to lifting an eigenclass of
\(\mathcal Q\) to a class in \(\mathcal B\) satisfying that same
generator equation. It is not CW4's map \(\delta^1:Q\to J(-1)\). CW4
constructs a different, geometric cohomology lift and proves its
boundary zero. Here the explicit nonzero boundary (GCF3.3) specifies its
own different lifting problem; it does not undo the geometric lift or
imply an obstruction at \(\tau\).

For the shifted normal action \(L+1\), exactly the same comparison has
entry \(F_0(\lambda-1)\). On \(\Re\lambda\in(0,1)\) this entry is
nonzero, since all zeros of \(F_0\) have real part in \((0,1)\). This
recovers the already proved separation of the \(Q\) and \(Q(-1)\)
spectra. It does not replace the full target \(J(-1)\) by its quotient.

\subsection{GCF4. The complete divisor module of the original
function}\label{gcf4.-the-complete-divisor-module-of-the-original-function}

Let \(\mathcal M\) be the sheaf of meromorphic functions on the
parameter plane. Write \([a]\) for the divisor at the scalar point
\(a\), and set \[
T=\sum_{n\ge1}[-2n],\qquad D=[1]-T.
\tag{GCF4.1}
\] The sum is locally finite, so its orders define a divisor on the
whole plane. The sheaf \(\mathcal O(D)\subset\mathcal M\) consists of
meromorphic functions \(g\) for which
\(\operatorname{ord}_a(g)+D(a)\ge0\) at every point. Locally it is a
free rank-one \(\mathcal O\)-module; at \(1\) its frame is
\((s-1)^{-1}\), and at \(-2n\) its frame is \(s+2n\).

The full factor \(c\) in (GCF0.3) has divisor exactly \(D\). Indeed
Gamma has no zeros, its simple pole at zero cancels the displayed factor
\(s\), the factor \(s-1\) gives a simple zero at \(1\), and its
remaining simple poles occur at every \(-2n\). The recurrence for Gamma
retains their residues in GCF5 below. The power of \(\pi\) is an
everywhere nonzero entire factor. There are no further zeros or poles.
Hence multiplication gives a global module isomorphism \[
\mu_c:\mathcal O(D)\xrightarrow{\sim}\mathcal O,
\qquad g\longmapsto cg,
\qquad f\longmapsto f/c.
\tag{GCF4.2}
\] The order inequalities prove both directions at every exceptional
point, so (GCF4.2) is not division by a supposed holomorphic unit there.

The original meromorphic \(\zeta\) belongs to \(\mathcal O(D)\): it has
its simple pole at \(1\), its simple zeros at every \(-2n\), and no
other poles. The complete comparison is \[
\begin{array}{ccc}
\mathcal O&\xrightarrow{\ \cdot\zeta\ }&\mathcal O(D)\subset\mathcal M\\
\mathrm{id}\downarrow&&\downarrow\mu_c\ \cong\\
\mathcal O&\xrightarrow{\ \cdot F_0\ }&\mathcal O.
\end{array}
\tag{GCF4.3}
\] It is a commutative diagram of actual holomorphic line modules with
their specified meromorphic embeddings. Its inverse recovers
\(\zeta=F_0/c\) at every point. Retaining \(D\), \(\mu_c\), and the
inclusion into \(\mathcal M\) retains the original zero and pole orders.
Keeping only the bottom scalar function would not specify those
embeddings.

The cokernel of either horizontal module map is supported precisely on
the nontrivial zero divisor, with length \(m_\rho\) at \(\rho\). To
prove this, at a nontrivial zero the multiplier is \(h^{m_\rho}\) times
its retained holomorphic unit, where \(h=s-\rho\); the quotient has
basis \(1,h,\ldots,h^{m_\rho-1}\). At every other point the horizontal
map is multiplication by a unit in the displayed local frames. GCF5
computes those frames and units at all cancellation points.

\subsection{GCF5. Every exceptional fibre and every infinitesimal
order}\label{gcf5.-every-exceptional-fibre-and-every-infinitesimal-order}

At \(s_0=-2n\), \(n\ge1\), put \(h=s+2n\). The exact local
factorizations are \[
\zeta(-2n+h)=h\,a_n(h),\qquad
c(-2n+h)=h^{-1}b_n(h),\qquad
F_0(-2n+h)=a_n(h)b_n(h),
\tag{GCF5.1}
\] where the holomorphic units and their constants are \[
a_n(0)=\zeta'(-2n),\qquad
b_n(0)=\frac{(-2n)(-2n-1)}8\pi^n\frac{2(-1)^n}{n!}
=\frac{n(2n+1)(-1)^n\pi^n}{2n!}.
\tag{GCF5.2}
\] Gamma's residue follows by applying \(\Gamma(z+1)=z\Gamma(z)\)
repeatedly to \(\Gamma(h/2)\). Nonvanishing of \(a_n(0)\) follows from
the original functional equation and the positive Dirichlet value at
\(1+2n\); explicitly \[
a_n(0)b_n(0)=
\frac{(1+2n)(2n)}8\pi^{-(1+2n)/2}
\Gamma((1+2n)/2)\zeta(1+2n)\ne0.
\tag{GCF5.3}
\] The fibre of the original target line is
\(h\mathcal O/h^2\mathcal O\), and the map \(\mu_c\) sends \(h a_n\) to
\(b_na_n\). It therefore retains \(\zeta'(-2n)\) and its full
multiplier, even though the bottom function is nonzero there.

At \(s=1\), put \(h=s-1\). Then \[
\zeta(1+h)=h^{-1}a_1^{\rm pole}(h),\quad
c(1+h)=h b_1^{\rm pole}(h),\quad
a_1^{\rm pole}(0)=1,\quad b_1^{\rm pole}(0)=1/8.
\tag{GCF5.4}
\] The original target fibre is \(h^{-1}\mathcal O/\mathcal O\). Its
residue coefficient \(1\) is carried to \(1/8\). At \(s=0\), the target
is \(\mathcal O\), \(c(0)=-1/4\), and \(\zeta(0)=-1/2\), again giving
\(F_0(0)=1/8\). Neither zero nor pole data has been silently removed.

All higher orders are equally explicit. At any point, write
\(\mathcal O(D)=h^{-d}\mathcal O\) and \(c=h^d b(h)\), with \(d=D(s_0)\)
and \(b(0)\ne0\). On the full length-\(r\) jet module, (GCF4.2) is \[
h^{-d}\mathcal O/h^{r-d}\mathcal O\longrightarrow
\mathcal O/h^r\mathcal O,
\qquad [h^{-d}a(h)]\longmapsto[b(h)a(h)],
\tag{GCF5.5}
\] for every \(r\ge1\). In ordered bases \(h^{-d},\ldots,h^{r-1-d}\) and
\(1,\ldots,h^{r-1}\), its matrix has entries \[
B_{ij}=\begin{cases}b^{(i-j)}(0)/(i-j)!,&i\ge j,\\0,&i<j,\end{cases}
\quad 0\le i,j<r,\qquad \det B=b(0)^r.
\tag{GCF5.6}
\] Its inverse is the same matrix rule with every coefficient of
\(1/b\). Thus every infinitesimal coefficient is retained; the
comparison is not only a statement about values or a first-order
approximation.

\subsection{GCF6. The transported involution on the original
target}\label{gcf6.-the-transported-involution-on-the-original-target}

The source functional equation gives \(F_0(s)=F_0(1-s)\). On the
original target module, the exact transported reflection is \[
(\mathscr R_Dg)(s)=\frac{c(1-s)}{c(s)}g(1-s),
\qquad \mu_c\mathscr R_D=\mathscr R\mu_c,
\quad (\mathscr Rf)(s)=f(1-s).
\tag{GCF6.1}
\] It is defined globally as a map of line modules over the reflection,
including exceptional points: by (GCF4.2), it is
\(\mu_c^{-1}\mathscr R\mu_c\). Thus its square is the identity, and
\(\mathscr R_D\zeta=\zeta\). The ratio is explicitly \[
\frac{c(1-s)}{c(s)}=
\frac{\frac{(1-s)((1-s)-1)}8\pi^{-(1-s)/2}\Gamma((1-s)/2)}
{\frac{s(s-1)}8\pi^{-s/2}\Gamma(s/2)}.
\tag{GCF6.2}
\] This is an isomorphism of the indicated line modules over reflection,
not an assertion that the ratio is an everywhere holomorphic unit on
scalar functions. In particular its zeros and poles implement the frames
in GCF5. On coefficient functions,
\(\mathscr R(L-\lambda)=-(L-(1-\lambda))\mathscr R\). On length-\(r\)
jets, reflection sends \(h^j\) to \((-1)^jh^j\) between the two centres.
Consequently the full generator comparison respects the reflected
parameter and every derivative sign. No radius or distance from \(\tau\)
was used to construct this action.

\subsection{GCF7. The original logarithmic derivative remains
complete}\label{gcf7.-the-original-logarithmic-derivative-remains-complete}

Off the zeros and poles, logarithmic differentiation of the actual
module comparison gives \[
\frac{F_0'}{F_0}
=\frac{\zeta'}{\zeta}
+\frac1s+\frac1{s-1}-\frac12\log\pi
+\frac12\frac{\Gamma'(s/2)}{\Gamma(s/2)}.
\tag{GCF7.1}
\] It then holds meromorphically. At a nontrivial zero its residue is
the complete multiplicity. At \(1\), the residues \(-1\) from
\(\zeta'/\zeta\) and \(+1\) from \(1/(s-1)\) cancel. At each \(-2n\),
the residues \(+1\) from the trivial zero and \(-1\) from Gamma cancel.
At zero the residues \(+1\) from \(1/s\) and \(-1\) from Gamma cancel;
\(\zeta\) is nonzero there. These are the same contributions encoded by
the frames in GCF5, with the power of \(\pi\) and every regular
derivative term still present in (GCF7.1).

\subsection{GCF8. Consequence for the requested weight
calculation}\label{gcf8.-consequence-for-the-requested-weight-calculation}

The actual coefficient inclusion has now produced a holomorphic
one-by-one determinant map, and (GCF4.3) returns it to the original
\(\zeta\) as an isomorphism of explicitly embedded divisor modules. The
precise spectral lifting boundary is (GCF3.2). It is nonzero at every
actual nontrivial zero, on or off the critical line, and retains
multiplicity through the complete jet modules. Thus it cannot by itself
distinguish the two locations. This is a calculation of the actual
boundary, rather than a hypothetical counterexample constructed outside
the programme.

CW4's geometric lifting boundary remains zero. Deligne's argument
additionally separates the weights of its entire source and target; no
such bound follows merely from the determinant map just derived. The
original full target \(J(-1)\), the degree factor, and the CW quotient
maps remain present. The new comparison supplies the exact original-zeta
object to which any further weight or trace argument must apply. It does
not replace that argument with a desired assumption.

Human source comparisons: Connes and Consani, \emph{Schemes over F1 and
zeta functions},
\href{https://arxiv.org/abs/0903.2024}{arXiv:0903.2024v3}, §5, for the
source summation/sheaf construction; Connes, \emph{The Riemann
Hypothesis: Past, Present and a Letter Through Time},
\href{https://arxiv.org/abs/2602.04022}{arXiv:2602.04022v1},
\texttt{rhready.tex} lines526--545, for the retained functional
equation; Deligne, \emph{La conjecture de Weil. II},
\href{https://www.numdam.org/item/PMIHES_1980__52__137_0/}{IHÉS52
(1980),137--252}, §3.6, for the distinct weight-separation lifting
mechanism. Source-reading coverage and hashes are recorded privately;
these citations do not assert a new reading of all external
dependencies.
